\section{Iteration and completion}\label{sec:completion}

We finish the proof of Theorem~\ref{thm:main} by iterating
Proposition~\ref{prop:one-round}.  The potential keeps a fixed positive
amount of weight at each remaining vertex, while the maximum remaining
degree decreases to a smaller order than the number of available
positive-weight entries.

\begin{proof}[Proof of Theorem~\ref{thm:main}]
Start from the empty partial coloring and the weights in
\eqref{eq:initial-state}, and set
\begin{equation}\label{eq:number-of-rounds}
 T_* = \left\lceil2\tau^{-1}\log\Delta\right\rceil.
\end{equation}
Perform $T_*$ rounds, stopping earlier if every vertex receives a
color.  We verify inductively that the mass assumption
\eqref{eq:mass-window} remains valid.  After $j\le T_*$ completed
rounds, every remaining vertex satisfies
\begin{equation}\label{eq:iteration-potentials}
 W_v\le K+j\xi,\qquad
 \gamma\log\Delta\,(K-W_v)\le F_v
 \le\frac{K^2}{C}\log\Delta+j\xi.
\end{equation}
The upper bounds follow from \eqref{eq:initial-state} and
\eqref{eq:round-conclusions}; the lower bound on $F_v$ follows from
the nonnegative entropy and edge terms in
\eqref{eq:vertex-potentials}.  Since
$T_*\xi=O(\Delta^{-0.04}\log\Delta)=o(1)$ and
$C=100K/\gamma$, these inequalities give
\begin{equation}\label{eq:iteration-mass-window}
 K-\frac K{100}-\frac{T_*\xi}{\gamma\log\Delta}
 \le W_v\le K+T_*\xi.
\end{equation}
For sufficiently large $\Delta$, both bounds lie inside $[K/2,2K]$.
Availability and \eqref{eq:weight-bounds} are preserved by the
proposition, so each successive round is justified.

If vertices remain after $T_*$ rounds, repeated use of the degree
inequality in \eqref{eq:round-conclusions} gives
\begin{equation}\label{eq:final-degree}
 \frac{\deg_{\mathrm{remaining}}(v)}{\Delta}
 \le(1-\tau)^{T_*}+\xi\sum_{j=0}^{T_*-1}(1-\tau)^j
 \le\Delta^{-2}+\frac{\xi}{\tau}
 =\Delta^{-2}+\Delta^{-0.04}.
\end{equation}
On the other hand, the mass lower bound and the upper bound on each
positive entry imply
\begin{equation}\label{eq:available-entry-count}
 \bigl|\{i\in L(v):w_i>0\}\bigr|
 \ge\frac K2\Delta^{0.99}.
\end{equation}
For large $\Delta$ this exceeds the remaining maximum degree, which
is at most $\Delta^{-1}+\Delta^{0.96}$ by \eqref{eq:final-degree}.
Color the remaining vertices in any order, choosing at each step a
positive-weight entry compatible with the colors already chosen.
Every such entry is available with respect to the partial coloring
from the rounds.  Each neighbor colored during this final procedure
forbids at most one entry, by the matching condition.  Inequality
\eqref{eq:available-entry-count} therefore guarantees a choice at
every step.

The resulting coloring chooses one entry from each restricted list
and avoids every forbidden pair, so it is also a coloring for the
original lists and matchings.  All lower thresholds on $\Delta$ in
the proof depend only on $r$.  Choosing a single
$\Delta_r\ge3$ above them, together with $C_r=C$ from
\eqref{eq:round-constants}, proves Theorem~\ref{thm:main}.
\end{proof}
